% META: {"id": "TEXP-COMPLEXES-TRIGO-POLYNOMES-QCM", "chapitre": "TEXP-COMPLEXES-TRIGO-POLYNOMES", "type_objet": "qcm", "genere_depuis": "chapitres/TEXP-COMPLEXES-TRIGO-POLYNOMES/qcm/TEXP-COMPLEXES-TRIGO-POLYNOMES-QCM.json", "status": "generated"}
% Fichier genere par scripts/build_qcm_tex.py — ne pas editer a la main.

\section*{\textcolor{chapcolor}{\MakeUppercase{Complexes : trigonométrie et polynômes}}}

\begin{center}
\textit{Pour chaque question, une seule réponse est exacte.}
\end{center}

\begin{enumerate}

\item[] \textbf{Capacité C1}

\item \textbf{[Q1]} La forme exponentielle de $z = 1 + i\sqrt{3}$ est :
  \begin{enumerate}[label=\Alph*.]
    \item $2e^{i\pi/3}$
    \item $2e^{i\pi/6}$
    \item $e^{i\pi/3}$
    \item $2e^{-i\pi/3}$
  \end{enumerate}

\item[] \textbf{Capacité C2}

\item \textbf{[Q2]} D'après la formule de Euler, $\cos(\theta) =$
  \begin{enumerate}[label=\Alph*.]
    \item $\frac{e^{i\theta} - e^{-i\theta}}{2}$
    \item $\frac{e^{i\theta} + e^{-i\theta}}{2i}$
    \item $\frac{e^{i\theta} + e^{-i\theta}}{2}$
    \item $e^{i\theta} + e^{-i\theta}$
  \end{enumerate}

\item \textbf{[Q3]} La formule de Moivre donne $(\cos\theta + i\sin\theta)^n =$
  \begin{enumerate}[label=\Alph*.]
    \item $\cos^n\theta + i\sin^n\theta$
    \item $\cos(n\theta) + i\sin(n\theta)$
    \item $n\cos\theta + in\sin\theta$
    \item $\cos(n\theta) - i\sin(n\theta)$
  \end{enumerate}

\item[] \textbf{Capacité C7}

\item \textbf{[Q4]} Les racines n-iemes de l'unité sont au nombre de :
  \begin{enumerate}[label=\Alph*.]
    \item 1
    \item 2
    \item n-1
    \item n
  \end{enumerate}

\item[] \textbf{Capacité C5}

\item \textbf{[Q5]} Si $P$ est un polynôme complexe et $P(a)=0$, alors :
  \begin{enumerate}[label=\Alph*.]
    \item $z-a$ divise $P(z)$
    \item $z+a$ divise toujours $P(z)$
    \item $P$ est le polynôme nul
    \item $a$ est l'unique racine de $P$
  \end{enumerate}

\end{enumerate}

% NEXUS-QCM-TEACHER-ONLY-BEGIN
\ifnxVersionProfesseur
\clearpage
\section*{Clé de correction — réservée au professeur}

\begin{center}
\begin{tabular}{lll}
\hline
Question & Capacité & Réponse exacte \\
\hline
Q1 & C1 & \textbf{A} \\
Q2 & C2 & \textbf{C} \\
Q3 & C2 & \textbf{B} \\
Q4 & C7 & \textbf{D} \\
Q5 & C5 & \textbf{A} \\
\hline
\end{tabular}
\end{center}

\subsection*{Diagnostic des réponses erronées}

\begin{itemize}
  \item \textbf{Q1}
  \begin{itemize}
    \item \textbf{B} — Erreur d'angle pi/6 au lieu de pi/3. Le module de 1+i racine(3) vaut 2 et son argument principal vaut pi/3, donc z=2 exp(i pi/3). \emph{Renvoi : C1.}
    \item \textbf{C} — Module faux (1 au lieu de 2). Le module de 1+i racine(3) vaut 2 et son argument principal vaut pi/3, donc z=2 exp(i pi/3). \emph{Renvoi : C1.}
    \item \textbf{D} — Erreur de signe d'argument. Le module de 1+i racine(3) vaut 2 et son argument principal vaut pi/3, donc z=2 exp(i pi/3). \emph{Renvoi : C1.}
  \end{itemize}
  \item \textbf{Q2}
  \begin{itemize}
    \item \textbf{A} — Formule de sin(theta) sans le i. En additionnant les formules d'Euler, cos(theta)=(exp(i theta)+exp(-i theta))/2. \emph{Renvoi : C2.}
    \item \textbf{B} — Facteur 2i au lieu de 2. En additionnant les formules d'Euler, cos(theta)=(exp(i theta)+exp(-i theta))/2. \emph{Renvoi : C2.}
    \item \textbf{D} — Oubli de diviser par 2. En additionnant les formules d'Euler, cos(theta)=(exp(i theta)+exp(-i theta))/2. \emph{Renvoi : C2.}
  \end{itemize}
  \item \textbf{Q3}
  \begin{itemize}
    \item \textbf{A} — Puissance appliquee aux fonctions. La formule de Moivre donne (cos theta+i sin theta)\textasciicircum{}n=cos(n theta)+i sin(n theta). \emph{Renvoi : C2.}
    \item \textbf{C} — Facteur n devant cos et sin. La formule de Moivre donne (cos theta+i sin theta)\textasciicircum{}n=cos(n theta)+i sin(n theta). \emph{Renvoi : C2.}
    \item \textbf{D} — Erreur de signe. La formule de Moivre donne (cos theta+i sin theta)\textasciicircum{}n=cos(n theta)+i sin(n theta). \emph{Renvoi : C2.}
  \end{itemize}
  \item \textbf{Q4}
  \begin{itemize}
    \item \textbf{A} — Seule la racine 1 comptee dans R. L'équation z\textasciicircum{}n=1 possède exactement les n solutions exp(2ikpi/n), pour k de 0 à n-1. \emph{Renvoi : C7.}
    \item \textbf{B} — Racines réelles uniquement. L'équation z\textasciicircum{}n=1 possède exactement les n solutions exp(2ikpi/n), pour k de 0 à n-1. \emph{Renvoi : C7.}
    \item \textbf{C} — Compte incorrect. L'équation z\textasciicircum{}n=1 possède exactement les n solutions exp(2ikpi/n), pour k de 0 à n-1. \emph{Renvoi : C7.}
  \end{itemize}
  \item \textbf{Q5}
  \begin{itemize}
    \item \textbf{B} — La racine $a$ fournit le facteur $z-a$ ; le facteur $z+a$ correspondrait à la racine $-a$. Le théorème de factorisation affirme que P(a)=0 si et seulement si z-a divise P(z). \emph{Renvoi : C5.}
    \item \textbf{C} — Un polynôme non nul peut avoir des racines ; l'annulation en un point n'implique pas l'annulation identique. Le théorème de factorisation affirme que P(a)=0 si et seulement si z-a divise P(z). \emph{Renvoi : C5.}
    \item \textbf{D} — La factorisation par $z-a$ n'interdit pas d'autres facteurs et donc d'autres racines. Le théorème de factorisation affirme que P(a)=0 si et seulement si z-a divise P(z). \emph{Renvoi : C5.}
  \end{itemize}
\end{itemize}
\fi
% NEXUS-QCM-TEACHER-ONLY-END
